\documentclass[a4paper,12pt]{article}
\usepackage[portuguese]{babel}
\usepackage{graphicx}
\usepackage{hyperref}

\title{Relatório de Projeto EDA - Primeira Fase}
\author{Bruno Paiva \\ a31496@alunos.ipca.pt}
\date{\today}

\begin{document}

\maketitle
\newpage

\section{Objetivo do Trabalho}
Este projeto teve como objetivos principais:
\begin{itemize}
    \item Praticar o uso de listas ligadas em C
    \item Aprender a ler dados de arquivos
    \item Organizar o código em funções bem definidas
    \item Documentar o projeto usando Doxygen
\end{itemize}

\section{Descrição do Problema}
O projeto simula uma cidade com várias antenas de transmissão, onde:
\begin{itemize}
    \item Cada antena possui:
    \begin{itemize}
        \item Posição no mapa (coordenadas X e Y)
        \item Frequência de operação (representada por um caractere)
    \end{itemize}
    
    \item Quando duas antenas com a \textbf{mesma frequência} estão alinhadas:
    \begin{itemize}
        \item Geram \textbf{efeitos nefastos} em certas posições
        \item Esses efeitos ocorrem em pontos específicos do alinhamento
    \end{itemize}
\end{itemize}

\newpage

\section{Solução Implementada}
\subsection{Estrutura de Dados Principal}
\begin{verbatim}
typedef struct Antena {
    int x, y;           // Posição (coordenadas)
    char frequencia;    // Frequência (A-Z)
    struct Antena *next;// Próxima antena na lista
} Antena;
\end{verbatim}

\subsection{Principais Funcionalidades}
\begin{itemize}
    \item \textbf{Carregar antenas}: Lê um arquivo texto com o mapa de antenas
    \item \textbf{Inserir/Remover}: Gerencia antenas individualmente
    \item \textbf{Detectar efeitos nefastos}: Calcula automaticamente as zonas com efeito nefasto
    \item \textbf{Exibir resultados}: Mostra uma tabela com:
    \begin{itemize}
        \item Todas as antenas
        \item Locais com efeitos nefastos
    \end{itemize}
\end{itemize}

\section{Exemplo de Funcionamento}
Para este mapa:
\begin{verbatim}
    .....
    .A.B.
    ..A..
    .....
\end{verbatim}

O programa detecta:
\begin{itemize}
    \item 2 antenas na frequência 'A'
    \item 2 locais com efeito nefasto entre elas
    \item 1 antena na frequência 'B' (sem efeitos)
\end{itemize}

\section{Conclusão}
Através deste projeto pude aprender e praticar:
\begin{itemize}
    \item Como criar e gerenciar listas ligadas
    \item Técnicas para processar arquivos
    \item Métodos para calcular posições em matrizes
    \item Boas práticas de documentação
\end{itemize}

\section{GitHub}
Link do repositório git

\url{https://github.com/yesisr/TP_EDA_Fase-1}

\end{document}